% Einheit 2 von 5 – Quader und Würfel (bank/koerper/e2.jsonl, zusammenbau v0.1)
%% TODO Zweigzeile Teil 1: was hier gelernt wird (Satz in Schülersprache aus dem Einheitstitel) – Einheit 2
\einheitenkopf[e2]{Einheit 2 von 5 · Quader und Würfel}
\zweigzeile{ab Kl. 5, je nach Buch bis Kl. 6 · P10 oft}
\setcounter{aufgabe}{13}

%% TODO Titel: Ich-kann-Satz fehlt in der Bank (Platzhalter: Kettenname) – Nr. 14
%% TODO Anweisung: ein Satz über der ersten Teilaufgabe fehlt in der Bank – Nr. 14
\begin{aufgabe}{Quader}
%% TODO \rechenplatz{n}: Schrittzahl des Grundfalls fehlt in der Bank (nur bei fester Schreibform, 2.3 b)
\begin{teile}
\teil Wie viel Wasser passt in das Aquarium? Gefragt ist das … \kreuz{Volumen} \kreuz{Oberfläche}
\teil Quader: Länge $6$ cm, Breite $3$ cm, Höhe $2$ cm – Volumen? \leerfeld[cm³]
\teil Ein Würfel hat $5$ cm Kantenlänge – Volumen? \leerfeld[cm³]
\teil $4\,500$ cm³ – wie viele Liter? \leerfeld[l]
\teil Quader: $7$ cm lang, $3$ cm breit, $2$ cm hoch – Oberfläche? \leerfeld[cm²]
\teil Ein würfelförmiger Hocker hat $5$ dm Kantenlänge – Oberfläche? \leerfeld dm²
\teil Quader: Volumen $90$ cm³, Länge $5$ cm, Breite $3$ cm – Höhe? \leerfeld[cm]
\teil Eine Betontreppe besteht aus zwei Quadern: unten $6$ dm × $4$ dm × $2$ dm, oben darauf $3$ dm × $4$ dm × $2$ dm. Volumen? \leerfeld dm³
\teil Ein Würfel hat die Kantenlänge $8$ cm. Gib sein Volumen an \hfill (P10 2026 FOR). \leerfeld[cm³]
\end{teile}
\end{aufgabe}

%% TODO Titel: Ich-kann-Satz fehlt in der Bank (Platzhalter: Kettenname) – Nr. 15
%% TODO Anweisung: ein Satz über der ersten Teilaufgabe fehlt in der Bank – Nr. 15
\begin{aufgabe}{Oberfläche ohne Deckel (Kiste, Aquarium)}
\begin{teile}
\teil Eine Holzkiste ohne Deckel ist $8$ dm lang, $5$ dm breit und $4$ dm hoch. Wie viel Holz braucht man? \leerfeld dm²
\end{teile}
\end{aufgabe}

%% TODO Titel: Ich-kann-Satz fehlt in der Bank (Platzhalter: Kettenname) – Nr. 16
%% TODO Anweisung: ein Satz über der ersten Teilaufgabe fehlt in der Bank – Nr. 16
\begin{aufgabe}{Würfelkante aus V (Kubikzahlen)}
\begin{teile}
\teil Würfel mit $27\,000$ cm³ Volumen – Kantenlänge? \leerfeld[cm]
\end{teile}
\end{aufgabe}

%% TODO Titel: Ich-kann-Satz fehlt in der Bank (Platzhalter: Kettenname) – Nr. 17
%% TODO Anweisung: ein Satz über der ersten Teilaufgabe fehlt in der Bank – Nr. 17
\begin{aufgabe}{Einheiten cm³, dm³ = l, m³ (1 l = 1000 cm³)}
\begin{teile}
\teil Ein Pool fasst $3$ m³ – wie viele Liter? \leerfeld[l]
\end{teile}
\end{aufgabe}

%% TODO Titel: Ich-kann-Satz fehlt in der Bank (Platzhalter: Kettenname) – Nr. 18
%% TODO Anweisung: ein Satz über der ersten Teilaufgabe fehlt in der Bank – Nr. 18
\begin{aufgabe}{Aquarium: Füllmenge in Litern, Füllhöhe}
\begin{teile}
\teil Ein Aquarium ist innen $50$ cm lang und $25$ cm breit. Es werden $50$ l Wasser eingefüllt. Wie hoch steht das Wasser? \leerfeld[cm]
\end{teile}
\end{aufgabe}

%% TODO Titel: Ich-kann-Satz fehlt in der Bank (Platzhalter: Kettenname) – Nr. 19
%% TODO Anweisung: ein Satz über der ersten Teilaufgabe fehlt in der Bank – Nr. 19
\begin{aufgabe}{Quader – Fehler finden, Begründen, Anwendung, Darstellungswechsel}
\begin{teile}
\teil Ein Quader ist $7$ cm lang, $2$ cm breit und $5$ cm hoch. Nils berechnet die Oberfläche: \rechnung{7 \cdot 2 \cdot 5 &= 70 \text{ cm}^2} Wo ist der Fehler?
\teil Warum ist das Volumen eines Quaders Länge mal Breite mal Höhe? Erkläre mit Würfeln von $1$ cm Kantenlänge.
\teil Ein Umzugskarton ist innen $62$ cm lang, $34$ cm breit und $36$ cm hoch. Im Katalog steht „ca. 75 Liter“. Stimmt die Angabe?
\teil Zeichne das Netz eines Quaders mit $7$ cm, $2$ cm und $1$ cm Kantenlänge und schreibe in jedes Rechteck seinen Flächeninhalt.

\rechenplatz{5}
\end{teile}
\end{aufgabe}
